\section{Лекция 1. Основы дискретной математики}

\subsection*{Примерная структура курса}

\begin{enumerate}
    \item Множества и отношения \hfill \(\sim 2\)
    \item Комбинаторика \hfill \(\sim 4\)
    \item Теория графов \hfill \(\sim 4\)
    \item Теория кодирования \hfill \(\sim 2\)
    \item Булевы функции \hfill \(\sim\) остальное
\end{enumerate}

\subsection{Множества (Наивная теория множеств)}

\begin{definition}
    \textdef{Множество} --- совокупность элементов.
\end{definition}

\begin{designation}
    \begin{align*}
        X \quad &\text{--- множества обозначаются заглавными буквами.} \\
        x \quad &\text{--- элементы множества обозначаются строчными (прописными) буквами.} \\
        x \in X \quad &\text{--- элемент } x \text{ принадлежит множеству } X. \\
        x \notin X \quad &\text{--- элемент } x \text{ не принадлежит множеству } X.
    \end{align*}
\end{designation}

\noindent Способы задания множеств:
\begin{itemize}
    \item Перечислением: \(\{x_1, x_2, \dots, x_n\}\);
    \item Характеристическим свойством: \(\{x \mid \varphi\}\) или \(\{x \in X \mid \varphi\}\).
\end{itemize}

\noindent \textbf{Специальные множества:}
\begin{itemize}
    \item Булево множество: \(\B = \{0, 1\}\);
    \item Числовые множества: \(\N, \Z, \Q, \R, \C\).
\end{itemize}

\subsubsection{Операции над множествами}

\begin{definition}
    Пусть \(X\) и \(Y\) --- некоторые множества. Определим следующие операции над ними:
    \begin{enumerate}
        \item \textdef{Объединение:} \(X \cup Y = \{x \mid x \in X \text{ или } x \in Y\}\);
        \item \textdef{Пересечение:} \(X \cap Y = \{x \mid x \in X \text{ и } x \in Y\}\);
        \item \textdef{Разность:} \(X \setminus Y = \{x \mid x \in X \text{ и } x \notin Y\}\);
        \item \textdef{Симметрическая разность:} \(X \triangle Y = (X \setminus Y) \cup (Y \setminus X) = (X \cup Y) \setminus  (Y \cap X)\);
        \item \textdef{Дополнение:} \(C_X Y = X \setminus Y\) (дополнение \(Y\) до \(X\), где \( Y \subseteq X \), то есть \(Y\) --- подмножество \(X\)).
    \end{enumerate}
\end{definition}

\subsubsection{Подмножества и равенство множеств}

\begin{definition}
    Множество \(X\) является \textdef{подмножеством} \(Y\) (\(X \subseteq Y\)), если имеет место импликация:
    \[
        (x \in X) \implies (x \in Y)
    \]
\end{definition}

\begin{definition}
    Множества \(X\) и \(Y\) называются \textdef{равными} (\(X = Y\)), если:
    \[
        X \subseteq Y \ \text{и} \ Y \subseteq X
    \]
\end{definition}

\begin{definition}
    \(X\) является \textdef{строгим подмножеством} \(Y\) (\(X \subset Y\)), если:
    \[
        X \subseteq Y \ \text{и} \ X \neq Y
    \]
\end{definition}

\begin{definition}
    \textdef{Пустое множество} \(\varnothing\) --- множество, не содержащее ни одного элемента.
\end{definition}

\begin{remark}
    Пустое множество является подмножеством любого множества (\(\forall X : \varnothing \subseteq X\)).
\end{remark}

\begin{remark}
    Следует различать пустое множество \(\varnothing\) и \(\{\varnothing\}\) --- множество, состоящее из
    пустого множества. Множество \(\varnothing\) не содержит элементов, а \(\{\varnothing\}\) содержит
    ровно один элемент --- само \(\varnothing\). Поэтому \(\varnothing \neq \{\varnothing\},\) хотя \(\varnothing \in \{\varnothing\}\) и \(\varnothing \subseteq \{\varnothing\}\).
\end{remark}


\subsubsection{Декартово произведение}

\begin{definition}
    \textdef{Совокупность} --- первичное, неопределяемое строго понятие, обозначающее
    мысленное объединение объектов в единое целое. 
\end{definition}

\begin{definition}
    \textdef{Кортеж} --- конечный упорядоченный набор из \(n\) элементов,
    обозначаемый \((x_1, x_2, \dots, x_n)\). В отличие от множества, в кортеже
    важен порядок, а элементы могут повторяться.
\end{definition}

\begin{definition}
    \textdef{Последовательность} --- функция, областью определения которой
    является множество натуральных чисел \(\N\) (бесконечная последовательность)
    или его начальный отрезок (конечная последовательность).
\end{definition}

\begin{definition}
    \textdef{Алфавит} --- конечное непустое множество, элементы которого
    называются символами или буквами.
\end{definition}

\begin{definition}
    \textdef{Слово конечной длины} (строка) --- конечная последовательность (кортеж)
    символов, выбранных из заданного алфавита.
\end{definition}

\begin{definition}
    \textdef{Пустое слово} --- слово, не содержащее ни одного символа
    (длина равна 0), обозначается \(\varepsilon\) или \(\Lambda\).
\end{definition}

\begin{definition}
    \textdef{Упорядоченная пара} \((x, y)\) --- совокупность, порядок элементов
    в которой важен.
\end{definition}

\begin{definition}
    \textdef{Конкатенация} (сцепление) слов --- бинарная операция, заключающаяся
    в приписывании одного слова в конец другого. Если даны два
    слова \(u = a_1 a_2 \dots a_k\) и \(v = b_1 b_2 \dots b_m\), то их конкатенацией
    называется слово \(uv = a_1 a_2 \dots a_k b_1 b_2 \dots b_m\) длины \(k + m\).
\end{definition}

\begin{remark}
    Свойства конкатенации и пустого слова:
    \begin{enumerate}
        \item Пустое слово является нейтральным элементом относительно
        операции конкатенации. Для любого слова \(w\) выполняется равенство: 
        \[w \varepsilon = \varepsilon w = w\]
        \item Операция конкатенации ассоциативна: 
        \[(uv)w = u(vw)\]
        \item В общем случае операция конкатенации некоммутативна: 
        \[uv \neq vu\]
    \end{enumerate}
\end{remark}

\begin{definition}
    \textdef{Декартово произведение} \(X \times Y = \{(x, y) \mid x \in X, y \in Y\}\).
\end{definition}

Обобщение на \(n\) множеств:
\[
    X_{1} \times X_{2} \times \dots \times X_{n}
    = \{\underbrace{(x_1, \dots, x_n)}_{\text{кортеж}} \mid x_{i} \in X_{i}, \ i = \overline{1, n}\},
    \quad n \in \N
\]

\begin{remark}
    Декартово произведение не ассоциативно:
    \[
        X_1 \times X_2 \times X_3 = \{(x_1, x_2, x_3)\} \quad \text{(множество троек)}
    \]
    \[
        (X_1 \times X_2) \times X_3 = \{((x_1, x_2), x_3)\} \quad \text{(множество пар, где первый элемент --- пара)}
    \]
\end{remark}

\(n\)-я декартова степень:
\[
    X^n = \underbrace{X \times X \times \dots \times X}_{n \text{ раз}}
        = \{\underbrace{(x_1, \dots, x_n)}_{\text{кортеж}} \mid x_i \in X, \ i = \overline{1, n}\},
        \quad n \in \N
\]

\begin{designation}
    \(\omega\) --- бесконечность.
\end{designation}

Счётная декартова степень:
\[
    X^\omega = \left\{
        \underbrace{(x_1, x_2, \dots, x_n, \dots)}_{
            \text{последовательность или } \omega \text{-слово}
        } \mid x_i \in X, \quad i = \overline{1, \infty}
    \right\}
\]

\begin{designation}
    \(X^{0}\) --- множество из одного элемента (пустого кортежа / пустого слова).
    \[
        X^0 = \{ () \} = \{ \varepsilon \}
    \]
\end{designation}

\begin{designation}
    Множество всех слов конечной длины (звезда Клини):
    \[
        X^* = \bigcup\limits_{k=0}^{\infty} X^k
    \]
    Множество всех слов конечной длины за исключением пустого слова:
    \[
        X^+ = \bigcup\limits_{k=1}^{\infty} X^k
    \]
    Множество всех слов (и конечные, и бесконечные):
    \[
        X^{\le \omega} = X^* \cup X^\omega
    \]
\end{designation}

\begin{definition}
    Множества $X$ и $Y$ называются \textdef{равномощными}, если между
    ними можно построить взаимно однозначное соответствие (биекцию). Обозначается \(X \sim Y\).
\end{definition}


\subsection{Мощность множеств}

\subsubsection{Счётные множества и кардинальные числа}

\begin{definition}
    \textdef{Мощность конечного множества} --- это число элементов в нём. Обозначается \(|X|\).
\end{definition}

\begin{definition}
    Множество называется \textdef{счётным}, если оно равномощно множеству натуральных чисел \(\N\).
\end{definition}

\begin{definition}
    \textdef{Кардинальное число} (или просто мощность) --- это характеристика множества, обобщающая понятие количества элементов на бесконечные множества. Для конечных множеств кардинальные числа совпадают с натуральными: \(0, 1, 2, \dots\). Для бесконечных множеств вводятся специальные символы (алефы).
\end{definition}

\begin{designation}
    Обозначения мощностей:
    \begin{align*}
        |\N| = |\Z| = |\Q| &= \aleph_0 \quad \text{(алеф-нуль)} \\
        |\R| &= \mathfrak{c} \quad \text{(континуум)}
    \end{align*}
    Шкала кардинальных чисел: \(0, 1, 2, \dots, \aleph_0, \aleph_1, \dots\)
\end{designation}

\begin{definition}
    Говорят, что множество \textdef{не более чем счётно}, если оно либо конечное, либо счётное.
\end{definition}

\begin{theorem}
    \emph{Свойства счётных множеств:}
    \begin{enumerate}
        \item Объединение конечного числа конечных множеств конечно.
        \item Объединение счётного числа конечных множеств не более чем счётно.
        \item Объединение конечного числа счётных множеств счётно.
        \item Объединение счётного числа счётных множеств счётно.
    \end{enumerate}
\end{theorem}

\vspace{1em}

\begin{proof}
    Докажем каждое из свойств.

    \textbf{1. Объединение конечного числа конечных множеств конечно.} \\
    Пусть даны конечные множества $A_1, A_2, \dots, A_n$. Так как каждое из них конечно, оно содержит конечное число элементов: $|A_i| = k_i \in \N \cup \{0\}$. Тогда мощность их объединения не превосходит суммы их мощностей: $\left|\bigcup_{i=1}^n A_i\right| \le \sum_{i=1}^n k_i$. Сумма конечного числа конечных чисел также конечна, следовательно, объединение является конечным множеством.

    \vspace{1em}
    \textbf{2. Объединение счётного числа конечных множеств не более чем счётно.} \\
    Пусть дано счётное семейство конечных множеств $A_1, A_2, A_3, \dots$. Так как каждое множество $A_i$ конечно, его элементы можно занумеровать: $A_i = \{a_{i1}, a_{i2}, \dots, a_{ik_i}\}$. Выпишем элементы всех множеств в одну последовательность, сначала элементы $A_1$, затем $A_2$, и так далее (при этом пропуская те элементы, которые уже встречались ранее). Поскольку каждый шаг добавляет лишь конечное число новых элементов, любой элемент из объединения рано или поздно получит свой натуральный номер. Следовательно, объединение можно занумеровать, и оно является не более чем счётным.

    \vspace{1em}
    \textbf{3. Объединение конечного числа счётных множеств счётно.} \\
    Пусть дано конечное число счётных множеств $A_1, A_2, \dots, A_n$. Так как каждое из них счётно, элементы любого $A_i$ можно выстроить в бесконечную последовательность: $A_i = \{a_{i1}, a_{i2}, a_{i3}, \dots\}$. Построим новую последовательность, беря по одному элементу из каждого множества по очереди («чередованием»): 
    \[a_{11}, a_{21}, \dots, a_{n1}, \ a_{12}, a_{22}, \dots, a_{n2}, \ \dots\]
    При такой нумерации элемент $a_{ij}$ получит позицию $(j-1)n + i$. Таким образом, мы установили сюръективное отображение из $\N$ в наше объединение. Так как объединение содержит хотя бы одно счётное подмножество, оно обязано быть бесконечным. Бесконечное не более чем счётное множество является счётным.

    \vspace{1em}
    \textbf{4. Объединение счётного числа счётных множеств счётно.} \\
    Пусть дано счётное число счётных множеств: $A_1, A_2, A_3, \dots$. 
    Так как каждое множество $A_i$ счётно, его элементы можно занумеровать: 
    $A_i = \{a_{i1}, a_{i2}, a_{i3}, \dots\}$.
    
    Расположим элементы всех множеств в виде бесконечной матрицы, где $i$-я строка содержит элементы множества $A_i$. Обойдём эту матрицу диагональной «змейкой» (метод Кантора), двигаясь по элементам, сумма индексов которых постоянна:
    
    \vspace{1em}
    \begin{minipage}[c]{0.55\textwidth}
    \begin{tikzpicture}[>={Stealth[scale=0.8]}]
    \matrix (m) [matrix of math nodes, row sep=1.5em, column sep=1.5em, nodes={anchor=center}]
    {
        A_1 & a_{11} & a_{12} & a_{13} & \dots \\
        A_2 & a_{21} & a_{22} & a_{23} & \dots \\
        A_3 & a_{31} & a_{32} & a_{33} & \dots \\
        \vdots & \vdots & \vdots & \vdots & \ddots \\
    };
    
    % Основной "зигзаг" Кантора по диагоналям
    \draw[->] (m-1-2) -- (m-1-3);
    \draw[->] (m-1-3) -- (m-2-2);
    \draw[->] (m-2-2) -- (m-3-2);
    \draw[->] (m-3-2) -- (m-2-3);
    \draw[->] (m-2-3) -- (m-1-4);
    
    % Переход на следующую диагональ и спуск вниз-влево
    \draw[->] (m-1-4) -- (m-1-5);
    \draw[->] (m-1-5) -- (m-2-4);
    \draw[->] (m-2-4) -- (m-3-3);
    
    % Дополнительная стрелка (переход к следующей диагонали)
    \draw[->] (m-3-3) -- (m-3-4);
    \end{tikzpicture}
    \end{minipage}
    \begin{minipage}[c]{0.4\textwidth}
    Порядок обхода: \\
    $a_{11}, a_{12}, a_{21}, a_{31}, a_{22}, a_{13}, a_{14}, a_{23}, a_{32}, \dots$ \\
    \\
    \textit{Если при обходе встречаются дубликаты (элемент принадлежит сразу нескольким множествам), мы их просто пропускаем.}
    \end{minipage}
    \vspace{1em}
    
    В результате каждый элемент из объединения $\bigcup A_i$ получит свой уникальный натуральный номер. Это устанавливает биекцию между объединением множеств и $\N$, следовательно, объединение счётно.
\end{proof}


\subsubsection{Метод диагонали Кантора}

\begin{theorem}[Кантора, 1-я формулировка]
    Множество $\B^\omega$ (множество всех бесконечных последовательностей из 0 и 1) не является счётным.
\end{theorem}

\begin{proof}
    Докажем от противного. Пусть множество \(\B^\omega\) счётно, тогда все его элементы
    можно занумеровать и выстроить в последовательность:
    \begin{align*}
        \alpha_{1} &= a_{11} a_{12} a_{13} \dots \\
        \alpha_{2} &= a_{21} a_{22} a_{23} \dots \\
        \alpha_{3} &= a_{31} a_{32} a_{33} \dots \\
        \vdots & \quad \vdots
    \end{align*}
    Рассмотрим элементы, стоящие на главной диагонали: \(\alpha = a_{11} a_{22} a_{33} \dots\)
    
    Построим новую последовательность \(\beta\), инвертируя каждый элемент диагонали: 
    \[\beta = (1-a_{11})(1-a_{22})(1-a_{33})\dots\]
    
    Очевидно, что \(\beta \in \B^\omega\), значит, она должна присутствовать в нашей пронумерованной таблице, то есть должен существовать номер \(n\) такой, что \(\alpha_{n} = \beta\). 
    
    Однако для любого \(n \in \N\), \(n\)-й символ последовательности \(\alpha_n\) равен \(a_{nn}\), в то время как \(n\)-й символ последовательности \(\beta\) равен \(1 - a_{nn} \neq a_{nn}\). Таким образом, \(\forall n \in \N: \alpha_{n} \neq \beta\). 
    Следовательно, последовательности \(\beta\) нет в нашем перечислении. Мы получили противоречие, значит, \(\B^\omega\) несчётно.
\end{proof}

\begin{remark}
    Рассмотренный способ доказательства называется \textdef{диагональным методом Кантора}.
\end{remark}

\subsubsection{Булеан и теорема Кантора}

\begin{definition}
    \textdef{Булеан} \(\mathcal{P}(X)\) (или \(2^X\)) --- множество всех
    подмножеств множества \(X\).
\end{definition}

\begin{example}
    \begin{align*}
        \mathcal{P}(\{0, 1, 2\}) &= \{\varnothing, \{0\}, \{1\},
            \{2\}, \{0, 1\}, \{0, 2\}, \{1, 2\}, \{0, 1, 2\}\} \\
        \mathcal{P}(\varnothing) &= \{\varnothing\} \quad (\text{при этом }
            |\varnothing| = 0, \ |\mathcal{P}(\varnothing)| = 1).
    \end{align*}
\end{example}

\begin{theorem}[Кантора, основная]
    Никакое множество не равномощно своему булеану:
    \[
        \forall X: X \nsim \mathcal{P}(X)
    \]
\end{theorem}

\begin{corollary}
    Мощность булеана всегда строго больше мощности самого множества \((|X| < |\mathcal{P}(X)|)\).
\end{corollary}

\begin{proof}
    Докажем от противного. Пусть \(\exists X: X \sim \mathcal{P}(X)\). Тогда существует
    взаимно однозначное соответствие (биекция) \(f: X \to \mathcal{P}(X)\).
    
    Рассмотрим множество \(Y = \{x \in X \mid x \notin f(x)\}\). Так как \(Y \subseteq X\), то \(Y \in \mathcal{P}(X)\). Поскольку \(f\) --- сюръекция (как часть биекции), должен существовать элемент \(y \in X\), такой что \(f(y) = Y\).
    
    Проверим, принадлежит ли элемент \(y\) множеству \(Y\):
    \begin{enumerate}
        \item Если \(y \in Y\), то по условию формирования множества \(Y\) имеем \(y \notin f(y)\). Но так как \(f(y) = Y\), получаем \(y \notin Y\). Противоречие.
        \item Если \(y \notin Y\), то \(y \notin f(y)\). Согласно правилу формирования множества \(Y\), элемент \(y\) должен быть включен в \(Y\), то есть \(y \in Y\). Снова противоречие.
    \end{enumerate}
    Оба возможных случая ведут к логическому парадоксу (подобному парадоксу Рассела). Следовательно, такого отображения \(f\) не существует, и \(X \nsim \mathcal{P}(X)\).
\end{proof}


\subsubsection{Континуум-гипотеза и бесконечные множества}

\noindent \textdef{Континуум-гипотеза} (Георга Кантора) --- это знаменитое математическое предположение о том, что между мощностью счётного множества (\(\aleph_0\)) и мощностью континуума (\(\mathfrak{c} = 2^{\aleph_0} = \aleph_1\)) не существует промежуточных мощностей. То есть, если мы берём любое бесконечное подмножество множества действительных чисел, то оно либо счётно, либо континуально. В стандартной аксиоматике теории множеств (ZFC) континуум-гипотеза является неразрешимой.

\begin{theorem}
    Любое бесконечное множество имеет счётное подмножество.
\end{theorem}

\begin{proof}
    Пусть $X$ --- произвольное бесконечное множество. Поскольку $X$ бесконечно, оно не пусто, значит, можно выбрать некий элемент $x_1 \in X$.
    
    Рассмотрим множество $X \setminus \{x_1\}$. Оно также не пусто, так как в противном случае множество $X$ состояло бы ровно из одного элемента и было бы конечным. Следовательно, мы можем выбрать элемент $x_2 \in X \setminus \{x_1\}$.
    
    Продолжая этот процесс по индукции, на $n$-м шаге мы рассматриваем разность $X \setminus \{x_1, x_2, \dots, x_{n-1}\}$. Это множество не пусто (иначе $X$ было бы конечным множеством из $n-1$ элементов), поэтому мы можем выбрать элемент $x_n \in X \setminus \{x_1, x_2, \dots, x_{n-1}\}$.
    
    Опираясь на аксиому выбора, мы можем продолжить этот процесс до бесконечности. В результате будет построена последовательность попарно различных элементов $x_1, x_2, \dots, x_n, \dots$, принадлежащих $X$. Множество этих элементов $A = \{x_1, x_2, \dots, x_n, \dots\}$ является подмножеством $X$ и, по построению, равномощно множеству натуральных чисел $\N$, то есть является счётным.
\end{proof}